\documentclass[aspectratio=169]{beamer}
% Shared Beamer configuration for MUS2640 lecture decks (included after \documentclass).
\usepackage[utf8]{inputenc}
\usepackage[T1]{fontenc}
\usepackage{lmodern}
\usepackage{graphicx}
\usepackage{booktabs}

\usetheme{Madrid}
\usecolortheme{dolphin}
\usefonttheme{professionalfonts}

\setbeamercovered{transparent}
\setbeamertemplate{navigation symbols}{}
\setbeamertemplate{footline}[frame number]

\newcommand{\coursename}{MUS2640 · Sensing Sound and Music}
\newcommand{\instituteuo}{University of Oslo · Department of Musicology / RITMO}


\title{Week 5 · Electroacoustics}
\subtitle{\coursename}
\institute{\instituteuo}
\date{}

\begin{document}

\begin{frame}
  \titlepage
\end{frame}

\begin{frame}{Aims}
  \begin{itemize}
    \item Trace the \textbf{signal chain}: sound $\rightarrow$ transducer $\rightarrow$ DSP/storage $\rightarrow$ reproduction.
    \item Terms: sampling, quantisation, latency; microphones \& loudspeakers at principle level.
    \item Connect fidelity constraints back to psychoacoustic thresholds.
  \end{itemize}
\end{frame}

\begin{frame}{Roadmap}
  \begin{enumerate}
    \item Analog vs digital; Nyquist intuition (why bandwidth limits matter).
    \item Microphones and pickups — polar patterns and placement (overview).
    \item Interfaces, monitors, headphones — colouration and critical listening.
    \item Effects \& dynamics as creative / corrective tools (categories).
  \end{enumerate}
\end{frame}

\begin{frame}{Why this week matters}
  \begin{itemize}
    \item Most students encounter music \textbf{mediated} — not only live in halls.
    \item Technology \textbf{selects} what can be heard in a mix or stream.
  \end{itemize}
\end{frame}

\begin{frame}{In-class}
  \begin{itemize}
    \item Whiteboard: draw a chain for a recording you make (phone, DAW, interface\ldots).
    \item Identify the top three places colouration or loss enters your chain.
  \end{itemize}
\end{frame}

\begin{frame}{Outside class}
  \begin{itemize}
    \item Sketch signal chain for one personal setup; chapter exercises.
  \end{itemize}
\end{frame}

\begin{frame}{Summary}
  \begin{itemize}
    \item Electroacoustics is where acoustics meets computation and listening habits.
    \item Next: \textbf{Time and rhythm} — metre, groove, micro-timing, entrainment.
  \end{itemize}
\end{frame}

\end{document}
